% id: PHY-0035
% title: ขนาดการกระจัดจากกราฟความเร็ว–เวลา
% camp: อื่น ๆ
% topic: การเคลื่อนที่แนวตรง
% year: 2557
% tags: IJSO, รอบที่ 1, การกระจัด, กราฟความเร็ว–เวลา
% origin: IJSO รอบที่ 1 2557 · ข้อ 35 ในภาพที่ผู้ใช้ส่ง
% review: ถอดจากภาพแล้ว · รอตรวจทาน
% diagram: PHY-0035.svg
% scan: sources/PHY-0035-original.png
% transcription-note: ต้นฉบับวงตัวเลือก 1 ไว้ แต่ไม่ได้บันทึกเป็นเฉลยที่ตรวจสอบแล้ว จุดต่ำสุดในกราฟไม่มีตัวเลขกำกับ จึงไม่เพิ่มป้ายค่าที่จุดนั้น
\documentclass[12pt]{article}
\usepackage{amsmath,amssymb,fontspec,polyglossia,tikz}
\setdefaultlanguage{thai}
\setmainfont{Noto Serif Thai}
\usetikzlibrary{arrows.meta}
\begin{document}
\textbf{ข้อ 35.} (IJSO รอบที่ 1 2557)
วัตถุชิ้นหนึ่งเคลื่อนที่ในหนึ่งมิติ ด้วยความเร็วที่เวลาต่าง ๆ ดังกราฟ โดยตอนแรกวัตถุมีความเร็ว $12\ \mathrm{m/s}$ และเมื่อเวลาผ่านไป $10\ \mathrm{s}$ วัตถุมีความเร็ว $18\ \mathrm{m/s}$ จงหาขนาดของการกระจัดของวัตถุในช่วงเวลา $10\ \mathrm{s}$ นี้

\begin{tikzpicture}[x=0.47897cm,y=0.11390cm,>=Stealth]
\draw[dashed,line width=.4pt] (0,18) -- (10,18);
\draw[dashed,line width=.4pt] (10,0) -- (10,18);
\draw[->,line width=.7pt] (-0.18,0) -- (13,0);
\draw[->,line width=.7pt] (0,-12) -- (0,23);
\draw[line width=.7pt] (4,-0.5853) -- (4,0.5853);
\node[anchor=center] at (4,-2.3411) {4};
\draw[line width=.7pt] (8,-0.5853) -- (8,0.5853);
\node[anchor=center] at (8,-2.3411) {8};
\draw[line width=.7pt] (10,-0.5853) -- (10,0.5853);
\node[anchor=center] at (10,-2.3411) {10};
\draw[line width=.7pt] (-0.1392,12) -- (0,12);
\node[anchor=east] at (-0.334,12) {12};
\draw[line width=.7pt] (-0.1392,18) -- (0,18);
\node[anchor=east] at (-0.334,18) {18};
\node[anchor=east] at (-0.2227,-1.8729) {0};
\node[anchor=west] at (13.2227,-0.2341) {$t$ (s)};
\node[anchor=west] at (-0.2227,24.8729) {$v$ (m/s)};
\draw[line width=.7pt] (0,12) -- (7,-9) -- (10,18);
\end{tikzpicture}

1. $24\ \mathrm{m}$

2. $36\ \mathrm{m}$

3. $48\ \mathrm{m}$

4. $60\ \mathrm{m}$
\end{document}
